\documentclass[12pt]{article}
\usepackage[utf8]{inputenc}
\usepackage[T1]{fontenc}
\usepackage[margin=1in]{geometry}
\usepackage{setspace}
\usepackage{xcolor}
\usepackage[colorlinks=true, urlcolor=blue!60!black, linkcolor=blue!60!black]{hyperref}
\setlength{\parindent}{0em}
\setlength{\parskip}{0.6em}
\renewcommand{\baselinestretch}{1.0}
\begin{document}
\begin{center}
{\LARGE \textit{Ex parte McCardle}, 74 U.S. (7 Wall.) 506 (1869)\par}
\vspace{0.5em}
{\large Case Brief --- Two-Page Edition\par}
\end{center}

\section{Memory Jogger}
Congress may repeal a statutory route of Supreme Court appellate jurisdiction even after argument; the Court dismissed McCardle's habeas appeal but noted that its 1789 habeas jurisdiction survived (\href{https://supreme.justia.com/cases/federal/us/74/506/}{74 U.S.\ 506}).

\section{Detailed Case Facts}
William McCardle, a Vicksburg newspaper editor who denounced Reconstruction, was held by the military for trial before a military commission. He sought habeas corpus under the Habeas Corpus Act of 1867, which authorized appeals to the Supreme Court.

\section{Procedural History}
The circuit court denied relief, and McCardle appealed. After the Supreme Court upheld its jurisdiction and heard argument, Congress, over President Johnson's veto, repealed the 1867 Act's appeal provision, including for pending appeals.

\section{Judicial Votes}
Dismissed, 8--0. Chief Justice Chase wrote for a unanimous Court.

\section{Holding}
The Exceptions Clause allows Congress to repeal appellate jurisdiction over pending cases, and courts may not inquire into legislative motive. Without jurisdiction, the Court must dismiss. The repeal reached only appeals under the 1867 Act, leaving jurisdiction under the Judiciary Act of 1789 intact.

\section{Analysis of Opinions}
Chase reasoned that appellate jurisdiction comes from the Constitution but subject to ``such Exceptions'' as Congress makes; the repeal was ``a plainer instance of positive exception'' than any. A repealed jurisdictional statute is treated as if it never existed, so the Court could only announce that fact and dismiss. His closing paragraph preserved the 1789 route, confirmed in \textit{Ex parte Yerger}. Later cases limited the principle: \textit{Klein} bars rules of decision disguised as jurisdiction strips, and \textit{Boumediene} requires an adequate substitute for habeas.

\section{Examples: Future Applications}
\textbf{Same-Side:} Congress may repeal direct appeals if certiorari remains available. \textbf{Opposite-Side:} stripping detainee habeas without an adequate substitute fails under \textit{Boumediene}. \textbf{Fence-Sitter:} eliminating every route of Supreme Court review in school-prayer cases.

\section{Critique}
On a broad reading, \textit{McCardle} gives Congress plenary power to take controversial subjects away from the Supreme Court; read narrowly, it allows closing one avenue while another remains. Henry Hart argued exceptions cannot destroy the Court's essential functions; critics find no such limit in the text.

\section{Key Quotations}
\begin{quote}
``Without jurisdiction, the court cannot proceed at all in any cause. Jurisdiction is power to declare the law, and, when it ceases to exist, the only function remaining to the court is that of announcing the fact and dismissing the cause.''
\end{quote}
\noindent --- Chief Justice Chase, Opinion of the Court, 74 U.S.\ at 514.

\end{document}
